\begin{statementbox}
	\textbf{\textcolor{CStatement}{条件}}
	\begin{description}[style=nextline, leftmargin=2.8em, labelsep=0.5em, itemsep=0.35em]
		\item[\textbf{\textcolor{CStatement}{条件一（定义域限制）：}}]
		      对数不等式要求 $x>0$；指数不等式对任意实数 $x$ 成立。
	\end{description}

	\medskip
	\textbf{\textcolor{CStatement}{结论}}
	\begin{description}[style=nextline, leftmargin=2.8em, labelsep=0.5em, itemsep=0.45em]
		\item[\textbf{\textcolor{CStatement}{结论一（对数切线）：}}]
		      \textbf{\textcolor{CStatement}{直观描述：}}
		      对数函数 $\ln x$ 的图象在点 $(1,0)$ 处的切线为 $y=x-1$，曲线位于切线下方。
		      \[
			      \boxed{\ln x\le x-1,\qquad x>0.}
		      \]

		\item[\textbf{\textcolor{CStatement}{结论二（指数切线）：}}]
		      \textbf{\textcolor{CStatement}{直观描述：}}
		      指数函数 $e^x$ 的图象在点 $(0,1)$ 处的切线为 $y=x+1$，曲线位于切线上方。
		      \[
			      \boxed{e^x\ge x+1,\qquad x\in\mathbb{R}.}
		      \]

		\item[\textbf{\textcolor{CStatement}{推广结论（常见变形）：}}]
		      \textbf{\textcolor{CStatement}{直观描述：}}
		      通过变量代换得到其他常用放缩形式。
		      \[
			      \boxed{\ln(1+t)\le t\quad (t>-1),}
			      \qquad
			      \boxed{e^x\ge e\,x\quad (x\in\mathbb{R}).}
		      \]
	\end{description}

	\medskip
	\textbf{\textcolor{CStatement}{等号/取等条件：}}
	\begin{itemize}[leftmargin=*, itemsep=0.5em]
		\item 对于 $\ln x\le x-1$，等号成立当且仅当 $x=1$；
		\item 对于 $e^x\ge x+1$，等号成立当且仅当 $x=0$；
		\item 对于 $\ln(1+t)\le t$，等号成立当且仅当 $t=0$；
		\item 对于 $e^x\ge e\,x$，等号成立当且仅当 $x=1$。
	\end{itemize}
\end{statementbox}